\documentclass[10pt,a4paper]{amsart}
\usepackage{amsmath,amssymb,amsthm,graphicx}
% For dealing with references we use the comment environment
\usepackage{verbatim}
\newenvironment{reference}{\comment}{\endcomment}
%\newenvironment{reference}{}{}
\newenvironment{slogan}{\comment}{\endcomment}
\newenvironment{history}{\comment}{\endcomment}

% For commutative diagrams we use Xy-pic
\usepackage[all]{xy}
% Theorem environments.
%
\theoremstyle{plain}
\newtheorem{theorem}[subsection]{Theorem}
\newtheorem{proposition}[subsection]{Proposition}
\newtheorem{lemma}[subsection]{Lemma}

\theoremstyle{definition}
\newtheorem{definition}[subsection]{Definition}
\newtheorem{example}[subsection]{Example}
\newtheorem{exercise}[subsection]{Exercise}
\newtheorem{situation}[subsection]{Situation}

\theoremstyle{remark}
\newtheorem{remark}[subsection]{Remark}
\newtheorem{remarks}[subsection]{Remarks}

\numberwithin{equation}{subsection}

% Macros
%
\def\lim{\mathop{\mathrm{lim}}\nolimits}
\def\colim{\mathop{\mathrm{colim}}\nolimits}
\def\Spec{\mathop{\mathrm{Spec}}}
\def\Hom{\mathop{\mathrm{Hom}}\nolimits}
\def\Ext{\mathop{\mathrm{Ext}}\nolimits}
\def\SheafHom{\mathop{\mathcal{H}\!\mathit{om}}\nolimits}
\def\SheafExt{\mathop{\mathcal{E}\!\mathit{xt}}\nolimits}
\def\Sch{\mathit{Sch}}
\def\Mor{\mathop{\mathrm{Mor}}\nolimits}
\def\Ob{\mathop{\mathrm{Ob}}\nolimits}
\def\Sh{\mathop{\mathit{Sh}}\nolimits}
\def\NL{\mathop{N\!L}\nolimits}
\def\CH{\mathop{\mathrm{CH}}\nolimits}
\def\proetale{{pro\text{-}\acute{e}tale}}
\def\etale{{\acute{e}tale}}
\def\QCoh{\mathit{QCoh}}
\def\Ker{\mathop{\mathrm{Ker}}}
\def\Im{\mathop{\mathrm{Im}}}
\def\Coker{\mathop{\mathrm{Coker}}}
\def\Coim{\mathop{\mathrm{Coim}}}

% Boxtimes
%
\DeclareMathSymbol{\boxtimes}{\mathbin}{AMSa}{"02}

%
% Macros for moduli stacks/spaces
%
\def\QCohstack{\mathcal{QC}\!\mathit{oh}}
\def\Cohstack{\mathcal{C}\!\mathit{oh}}
\usepackage{cleveref}
\crefname{originallemma}{original lemma}{original lemmas}
\Crefname{originallemma}{Original lemma}{Original lemmas}
\makeatletter
\newlabel{categories-lemma-yoneda}{{001P}{source}{}{sourcepointer-1}{}}
\newlabel{categories-lemma-yoneda@cref}{{[originallemma][0][]001P}{[1][0][]source}{}{}{}}
\makeatother
\crefname{originallemma}{original lemma}{original lemmas}
\Crefname{originallemma}{Original lemma}{Original lemmas}
\makeatletter
\newlabel{lemma-commutes-with-countable-sums}{{094A}{source}{}{sourcepointer-2}{}}
\newlabel{lemma-commutes-with-countable-sums@cref}{{[originallemma][0][]094A}{[1][0][]source}{}{}{}}
\makeatother
\crefname{originallemma}{original lemma}{original lemmas}
\Crefname{originallemma}{Original lemma}{Original lemmas}
\makeatletter
\newlabel{lemma-homological-functor-kernel}{{05RD}{source}{}{sourcepointer-3}{}}
\newlabel{lemma-homological-functor-kernel@cref}{{[originallemma][0][]05RD}{[1][0][]source}{}{}{}}
\makeatother
\title{Brown representability for compactly generated triangulated categories}
\author{The Stacks Project Authors}
\date{stacks-project:derived.tex:local-snapshot-20261011; GNU FDL 1.2 or later}
\begin{document}
\maketitle
\noindent\textit{Editorial note.} The compact-generator cellular approximation lemma is included with its complete proof, not granted for free and not counted separately.\par\medskip
\noindent\textit{Editorial note.} This is the compactly generated theorem. The broader brown-bis claim and its unresolved proof issue are not used. Original attribution to Neeman is retained.\par\medskip
\noindent\textit{Editorial scope.} One original target with its local construction proved in full. Classical category primitives are explicitly allowed.\par\medskip
\noindent\textit{Allowed standard prerequisites (editorial).}\begin{itemize}
\item Yoneda lemma and cohomological functors. An element of H(X) gives a natural map Hom(-,X)→H; exact triangles yield long exact sequences.
\item Compact objects and countable mapping telescopes. A compact C has Hom(C,hocolim X\_n)=colim Hom(C,X\_n). A compact generating family detects zero objects through all shifts.
\item Localizing closure of the isomorphism locus. The locus where a natural transformation of cohomological functors is an isomorphism on every shift is triangulated and closed under summands; if both preserve coproducts as products it is closed under coproducts.
\end{itemize}\medskip
\ifdefined\hypertarget\hypertarget{sourcepointer-1}{}\fi
\noindent\textit{Source reference.} Original lemma 001P: focused-seeds/ACL/sources/stacks-snapshot-20261011/categories.tex, source lines 540--558. Standard Yoneda lemma. Stable Stacks tag displayed.\par\medskip
\ifdefined\hypertarget\hypertarget{sourcepointer-2}{}\fi
\noindent\textit{Source reference.} Original lemma 094A: focused-seeds/ACL/sources/stacks-snapshot-20261011/derived.tex, source lines 10405--10422. Explicit compactness/telescope identity permitted as standard. Stable Stacks tag displayed.\par\medskip
\ifdefined\hypertarget\hypertarget{sourcepointer-3}{}\fi
\noindent\textit{Source reference.} Original lemma 05RD: focused-seeds/ACL/sources/stacks-snapshot-20261011/derived.tex, source lines 1812--1825. Standard long-exact-sequence closure verification, not representability. Stable Stacks tag displayed.\par\medskip
\begin{lemma}
\label{lemma-brown}
\begin{reference}
\cite[Theorem 3.1]{Neeman-Grothendieck}.
\end{reference}
Let $\mathcal{D}$ be a triangulated category with direct sums which is
compactly generated. Let $H : \mathcal{D} \to \textit{Ab}$ be a contravariant
cohomological functor which transforms direct sums into products.
Then $H$ is representable.
\end{lemma}

\begin{lemma}
\label{lemma-write-as-colimit}
Let $\mathcal{D}$ be a triangulated category with direct sums.
Let $E_i$, $i \in I$ be a family of compact objects of $\mathcal{D}$
such that $\bigoplus E_i$ generates $\mathcal{D}$.
Then every object $X$ of $\mathcal{D}$ can be written as
$$
X = \text{hocolim} X_n
$$
where $X_1$ is a direct sum of shifts of the $E_i$ and each transition
morphism fits into a distinguished triangle
$Y_n \to X_n \to X_{n + 1} \to Y_n[1]$
where $Y_n$ is a direct sum of shifts of the $E_i$.
\end{lemma}

\begin{proof}
Set $X_1 = \bigoplus_{(i, m, \varphi)} E_i[m]$ where the direct sum is over
all triples $(i, m, \varphi)$ such that $i \in I$, $m \in \mathbf{Z}$
and $\varphi : E_i[m] \to X$. Then $X_1$ comes equipped with a canonical
morphism $X_1 \to X$. Given $X_n \to X$ we set
$Y_n = \bigoplus_{(i, m, \varphi)} E_i[m]$ where the direct sum is over
all triples $(i, m, \varphi)$ such that $i \in I$, $m \in \mathbf{Z}$, and
$\varphi : E_i[m] \to X_n$ is a morphism such that $E_i[m] \to X_n \to X$
is zero. Choose a distinguished triangle
$Y_n \to X_n \to X_{n + 1} \to Y_n[1]$
and let $X_{n + 1} \to X$ be any morphism such that $X_n \to X_{n + 1} \to X$
is the given one; such a morphism exists by our choice of $Y_n$.
We obtain a morphism $\text{hocolim} X_n \to X$ by the construction
of our maps $X_n \to X$. Choose a distinguished triangle
$$
C \to \text{hocolim} X_n \to X \to C[1]
$$
Let $E_i[m] \to C$ be a morphism. Since $E_i$ is compact,
Lemma \ref{lemma-commutes-with-countable-sums}
implies the composition $E_i[m] \to C \to \text{hocolim} X_n$ factors through
$X_n$ for some $n$, say by $E_i[m] \to X_n$. Then the
construction of $Y_n$ shows that the composition
$E_i[m] \to X_n \to X_{n + 1}$ is zero. In other words, the composition
$E_i[m] \to C \to \text{hocolim} X_n$ is zero. This means that our
morphism $E_i[m] \to C$ comes from a morphism $E_i[m] \to X[-1]$.
The construction of $X_1$ then shows that such morphism lifts to
$\text{hocolim} X_n$ and we conclude that our morphism $E_i[m] \to C$
is zero. The assumption that $\bigoplus E_i$ generates $\mathcal{D}$
implies that $C$ is zero and the proof is done.
\end{proof}

\begin{proof}
Let $E_i$, $i \in I$ be a set of compact objects such that
$\bigoplus_{i \in I} E_i$ generates $\mathcal{D}$. We may and do assume
that the set of objects $\{E_i\}$ is preserved under shifts. Consider pairs
$(i, a)$ where $i \in I$ and $a \in H(E_i)$ and set
$$
X_1 = \bigoplus\nolimits_{(i, a)} E_i
$$
Since $H(X_1) = \prod_{(i, a)} H(E_i)$ we see that $(a)_{(i, a)}$
defines an element $a_1 \in H(X_1)$. Set $H_1 = \Hom_\mathcal{D}(- , X_1)$.
By Yoneda's lemma (Categories, Lemma \ref{categories-lemma-yoneda})
the element $a_1$ defines a natural transformation $H_1 \to H$.

\medskip\noindent
We are going to inductively construct $X_n$ and transformations
$a_n : H_n \to H$ where $H_n = \Hom_\mathcal{D}(-, X_n)$.
Namely, we apply the procedure
above to the functor $\Ker(H_n \to H)$ to get an object
$$
K_{n + 1} = \bigoplus\nolimits_{(i, k),\ k \in \Ker(H_n(E_i) \to H(E_i))} E_i
$$
and a transformation $\Hom_\mathcal{D}(-, K_{n + 1}) \to \Ker(H_n \to H)$.
By Yoneda's lemma the composition $\Hom_\mathcal{D}(-, K_{n + 1}) \to H_n$
gives a morphism $K_{n + 1} \to X_n$. We choose
a distinguished triangle
$$
K_{n + 1} \to X_n \to X_{n + 1} \to K_{n + 1}[1]
$$
in $\mathcal{D}$. The element $a_n \in H(X_n)$ maps to zero
in $H(K_{n + 1})$ by construction. Since $H$ is cohomological
we can lift it to an element $a_{n + 1} \in H(X_{n + 1})$.

\medskip\noindent
We claim that $X = \text{hocolim} X_n$ represents $H$. Applying $H$
to the defining distinguished triangle
$$
\bigoplus X_n \to
\bigoplus X_n \to X \to
\bigoplus X_n[1]
$$
we obtain an exact sequence
$$
\prod H(X_n) \leftarrow 
\prod H(X_n) \leftarrow 
H(X)
$$
Thus there exists an element $a \in H(X)$ mapping to $(a_n)$
in $\prod H(X_n)$. Hence a natural transformation
$\Hom_\mathcal{D}(- , X) \to H$ such that
$$
\Hom_\mathcal{D}(-, X_1) \to
\Hom_\mathcal{D}(-, X_2) \to
\Hom_\mathcal{D}(-, X_3) \to \ldots \to
\Hom_\mathcal{D}(-, X) \to H
$$
commutes. For each $i$ the map $\Hom_\mathcal{D}(E_i, X) \to H(E_i)$
is surjective, by construction of $X_1$. On the other hand, by construction
of $X_n \to X_{n + 1}$ the kernel of $\Hom_\mathcal{D}(E_i, X_n) \to H(E_i)$
is killed by the map
$\Hom_\mathcal{D}(E_i, X_n) \to \Hom_\mathcal{D}(E_i, X_{n + 1})$.
Since
$$
\Hom_\mathcal{D}(E_i, X) = \colim \Hom_\mathcal{D}(E_i, X_n)
$$
by Lemma \ref{lemma-commutes-with-countable-sums}
we see that $\Hom_\mathcal{D}(E_i, X) \to H(E_i)$ is injective.

\medskip\noindent
To finish the proof, consider the subcategory
$$
\mathcal{D}' =
\{Y \in \Ob(\mathcal{D}) \mid \Hom_\mathcal{D}(Y[n], X) \to H(Y[n])
\text{ is an isomorphism for all }n\}
$$
As $\Hom_\mathcal{D}(-, X) \to H$ is a transformation between
cohomological functors,
the subcategory $\mathcal{D}'$ is a strictly full, saturated, triangulated
subcategory of $\mathcal{D}$ (details omitted; see proof of
Lemma \ref{lemma-homological-functor-kernel}). Moreover, as both
$H$ and $\Hom_\mathcal{D}(-, X)$ transform direct sums into products,
we see that direct sums of objects of $\mathcal{D}'$ are in $\mathcal{D}'$.
Thus derived colimits of objects of $\mathcal{D}'$ are in $\mathcal{D}'$.
Since $\{E_i\}$ is preserved under shifts, we see that $E_i$
is an object of $\mathcal{D}'$ for all $i$. It follows from
Lemma \ref{lemma-write-as-colimit} that $\mathcal{D}' = \mathcal{D}$
and the proof is complete.
\end{proof}

\end{document}
